\section{Висновки}
\label{sec:conclusions}

\paragraph{Аналітичні результати.}
За ненасиченої двофакторної моделі з шумом нульового середнього і скінченної
дисперсії та зазначених припущень щодо шляху однофакторний показник МНК
задовольняє $\ahat_{m}\xrightarrow{\Prob}\gamma\lambda-\beta$
(Пропозиція~\ref{prop:identifiability}) --- зміщення через пропущену змінну.
Протилежні монотонні обмеження дають замкнене цілочислове вікно допустимості
(Пропозиція~\ref{prop:window}).

\paragraph{Експериментально підтримані спостереження.}
За \RNM{} на чотириядерній типовій CI $\beta=\resBeta\pm\resBetaSE$ і
$\gamma=\resGamma\pm\resGammaSE$ при $\q=0.25$; $\beta$ стабільний за
фіксованої квоти, але залежить від її рівня; повторне семплювання шляху
узгоджене з $\gamma\lambda-\beta$. Детектор близький до випадкового, тож вікно
визначає лише пропускна здатність. Подвоєний попит або WAN~$\beta$ потребує
$\kappaScale=\resKappaDoubled$ або $\approx\resKappaWan$ відповідно замість
$\resKappa$; ці значення спираються на неперевірену лінійність пропускної
здатності за квотою.
Числові результати специфічні для розгортання.

\paragraph{Інженерні наслідки.}
Внесок --- відтворювана методика ресурсно-нормованого калібрування та
планування потужності; компроміс «пропускна здатність --- детекція»
сформульовано, але з виміряним детектором не продемонстровано. Інші
розгортання потребують перекалібрування $(\beta,\gamma)$ і $(\pd,\pf,\rho)$ зі
збереженням того самого робочого процесу та публічних артефактів.

\subsection*{Конфлікт інтересів}
Автори заявляють про відсутність конфлікту інтересів.

\subsection*{Фінансування}
Дослідження виконано без фінансової підтримки.

\subsection*{Доступність даних}
Означення робочого процесу, генератори, код ін'єкції відмов і аналізу та сирі
артефакти доступні за адресою
\url{https://github.com/bpenyak/Resource_Normalized_Scaling_Laws_for_Byzantine_Consensus}.

\subsection*{Використання штучного інтелекту}
AI-допоміжні інструменти використано для мовного редагування та верстки. Усі
наукові результати, доведення, дизайн і аналіз виконані авторами.

\subsection*{Внесок авторів}
\textbf{Пеняк Б.\,О.}: концептуалізація, методологія, формальний аналіз,
програмне забезпечення, дослідження, курація даних, написання --- первинний
текст, візуалізація.
\textbf{Любінський Б.\,Б.}: наукове керівництво, валідація, написання ---
рецензування та редагування.
